% a04_Bh_bulk5  Proof of app:a04:lem:B2, part 5: the averaging part along the diagonal pairs, the Dirichlet weight, conclusion

\subsubsection*{B2.9 The averaging part along the diagonal pairs}

\begin{lemma}[gradient of the coarse-line pair difference]\label{app:a04:lem:Lam}
Let $g\defeq(\Gamma_t-\Gamma_{t/5})|_{\Lambda_t\times\Lambda_t}$ (both lines with the same $b$). Then
\[
|\nabla^{(t)}\Gamma_t-\nabla^{(t)}\Gamma_{t/5}|\le Ct^{1/2}(1+\ln t^{-1})^3,\qquad|\nabla^{(t)}\Gamma_{t/5}-\nabla^{(t/5)}\Gamma_{t/5}|\le Ct(1+\ln t^{-1})^3\quad\text{on }\Lambda_t\times\Lambda_t.
\]
\end{lemma}
\begin{proof}
$A\defeq\sup|g|\le Ct$ by ($\Gamma$4), and $B\defeq\sup|\nabla^{(t)}\nabla^{(t)}g|\le C(1+\ln t^{-1})^3$ by (G2) (coarse second differences of $\Gamma_{t/5}$ are averages of fine ones). Discrete Landau inequality: $g(u+Nte_\mu)-g(u)=Nt\,\nabla^{(t)}_\mu g(u)+t^2\sum_{i<N}\sum_{i'<i}\nabla\nabla g$, so $|\nabla^{(t)}g|\le2A/(Nt)+NtB/2$; take $N=\lfloor t^{-1/2}\rfloor$. The second bound: a coarse difference quotient of $\Gamma_{t/5}$ is the average of five fine ones along the step, and consecutive fine ones differ by $\le(t/5)B$.
\end{proof}

\begin{proposition}[the averaging part, increment]\label{app:a04:prop:AVb}
For $f=1$, all $k$ and $1\le j\le k-J_0$, with $q$ as in \cref{app:a04:prop:AVa} (enlarged if necessary) and $C$ uniform in $m',\mu'\in(0,m_0]$:
\[
\big|\bar{\mathfrak a}^{(j)}_k-\bar{\mathfrak a}^{(j+1)}_{k+1}\big|\le Ce^2\eta^{-1}\big[t^{1/2}(1+j)^q+5^{-2j}\big].
\]
\end{proposition}
\begin{proof}
Both members are evaluated at unit block scale on the same unit cell; the coarse member has spacing $t$, coupling $\mathfrak a_j$ and slices $i<j$, the finer one spacing $t/5$, coupling $\mathfrak a_{j+1}$ and slices $i'\le j$. Every term keeps a cycle through $x_0$ with unit-scale-decaying lines, so all vertex sums converge uniformly in $m',\mu'$ (the mass-free decay (G3)). Telescope one factor at a time in four steps.

\emph{(1) Unmatched slices} (slice $0$ of the coarser member, the two finest slices of the finer member): $O(t\cdot\mathrm{polylog})$ by (S5) and \cref{app:a04:lem:C}(i),(ii) ($\ell_0=t$ in both members).

\emph{(2) Common slices $i\ge1$ at coarse vertex points.} (S5) and \cref{app:a04:lem:C}(iv) give $Ct$ per slice on $\cV$ and $\cK$ (total $Ctj$). On $\cK_{\mathsf b}$: per slice $|\Delta\cK_{i,\mathsf b}|\le Ct/\ell_i$ (\cref{app:a04:lem:T}(T1)), which is not small for small $i$; in the region $|\mathsf b-x|\le r_i$ the E4 cost is $(t/\ell_i)\,r_i^2(1+r_i)R_t^3$ (the ball sum of $|G(x,\mathsf b_+)|\le C|x-\mathsf b|_t^{-1}+C$ over $|x-\mathsf b|\le r_i$ is $\le C(r_i^2+r_i^3)$), and $\sum_i(t/\ell_i)r_i^2(1+r_i)R_t^3\le Ct\sum_i\ell_i(i+1)^3R_t^3\le Ct(1+j)^3R_t^3$ ($r_i\le C\ell_i(i+1)$); in the shell $|\mathsf b-x|>r_i$ the pointwise bound carries no decay and would only give $\sum_it/\ell_i=O(1)$, so there \cref{app:a04:lem:T}(T3) is used, with the per-slice cost $Ct(1+j)^3R_t^3$ for both Wick families: the shell part is $\le Ct(1+j)^4R_t^3$.

\emph{(3) Lines at coarse points.} $C$-part: $|C_t-C_{t/5}|\le Ct^2|r|^{-3}$ and $|\nabla^{(t)}C_t-\nabla^{(t/5)}C_{t/5}|\le Ct|r|^{-3}$ (G1), summable against one more line at the cost of one logarithm; $\Gamma$-part: $|\Gamma_t-\Gamma_{t/5}|\le Ct+C5^{-2j}$ (G2); gradient-$\Gamma$ part: \cref{app:a04:lem:Lam}, $\le Ct^{1/2}(1+\ln t^{-1})^3$. This is the only step with rate $t^{1/2}$.

\emph{(4) Fine vertex points $\to$ coarse.} Kernel oscillation by \cref{app:a04:lem:C}(v) ($\le Ctj^2$ for $\cK'_{\mathsf b}$ in the region $|\mathsf b-x|\le r_i$); line oscillation $t(|r|_t^{-2}+1)$; gradient-line oscillation $Ct|r|_t^{-3}+Ct(1+\ln t^{-1})^3$ (G1), (G2); near-diagonal points ($|r|<2t$) bounded directly, $O(t)$. Shell: by \cref{app:a04:lem:C}(iii) the in-cell difference $D_{i'}(x',\mathsf b)\defeq\cK'_{i',\mathsf b}(x')-\cK'_{i',\mathsf b}(x)$ ($x'$ in the coarse cell of $x$) is the sum of $(t/5)S'_{i'}(\mathsf b',\mathsf b)$ over the $\le3(L-1)$ fine bonds $\mathsf b'$ of $\Gamma_{x,x'}$, all within $3t$ of $x$. With $\ell'\defeq5^{i'}t/5$ the scale of fine slice $i'$ ($0\le i'\le j$), $\lambda'\defeq2L\ell'/\kappa_0$, $\rho'\defeq(6/\kappa_0)L\ell'\ln(L^2\ell'/t)\le C\ell'(1+j)$ and $r'\defeq\rho'+3t$, (S2) and (S4) of the finer member give $|D_{i'}(x',\mathsf b)|\le C(t/\ell')\min(1,e^{-(r-r')/\lambda'})$, $r=|\mathsf b-x|$. The bond move ($\mathsf b\to$ the parallel bond based at the coarse point of its cell, $|\mathsf b-\mathsf b_c|\le3t$) obeys the same form $\min(Ct/\ell',CL^3e^{-(r-r')/\lambda'})$: pointwise by (S3) ($|\nabla S'_{i'}|\le C\ell'^{-2}$, $\le C_L\ell'L/t$ bonds of weight $t/5$, step $\le3t$), tail by (S4) applied to each of the two terms. \cref{app:a04:lem:T}(T3) for both Wick families ($t/\ell'\le L$; for the bond move with the split at $r'+ML\ell'$) gives per slice $\le Ct(1+j)^3$, over the $\le j+1$ fine slices $\le Ct(1+j)^4$, times the outer factor $CR_t^3$.

\emph{Coupling change.} $|\mathfrak a_j-\mathfrak a_{j+1}|\le C5^{-2j}$ with (G2) ($|\partial_b\Gamma|\le C$) and (S5): $C5^{-2j}$.

Summing (1)--(4): $|\bar{\mathfrak a}^{(j)}_k-\bar{\mathfrak a}^{(j+1)}_{k+1}|\le Ce^2\eta^{-1}[t^{1/2}(1+j)^q+5^{-2j}]$, the factors $e^2\eta^{-1}$ as in \cref{app:a04:prop:AVa}.
\end{proof}

\subsubsection*{B2.10 The averaging part with the Dirichlet weight, and the conclusion}

\begin{proposition}[the averaging part, weight $f_D$]\label{app:a04:prop:AVf}
\Cref{app:a04:prop:AVa,app:a04:prop:AVb} hold for the weight $f_D$ of \cref{app:a04:def:bulk}, with constants depending in addition on $\norm{f_D}_\infty$.
\end{proposition}
\begin{proof}
\emph{Setting.} The photon kernels are unchanged (the photon averaging is unchanged, \cite[p.~328]{K2}, so \cite[(2.17)]{K1} and \cref{app:a04:lem:S} hold verbatim). The vertex is $V^f_{\mathrm{av}}=-\frac b2\sum_y[|Q^f(A)\phi|^2-|Q^f\phi|^2]$ with $(Q^f(A)\phi)(y)=\sum_{x\in B(y)}t^3f(x)U(e'A(\Gamma_{y,x}))\phi(x)$ \cite[(3.2)]{K2}, the weight multiplying the fine field at every slice; $Q^f_{j+1}=Q_1Q^f_j$, $P^f_{i+1}\le P^f_i$ as before; the weight contributes bounded factors $\le\norm{f_D}_\infty^2$. \emph{Lines.} $G^f=C_t-\Gamma^f_t$ with the Dirichlet-weight clause of \cref{app:a04:lem:B0} (the same $\Gamma_0$, lower bound (A3) for $j\ge1$), \cref{app:a04:lem:CT} with $f_D$, and the gradient and second-difference bounds of \cref{app:a04:lem:G} with the D-aliases contributing $\le|p|^i\,Ct^2(1+\ln t^{-1})^3\le C\cdot\mathrm{polylog}$ for $i\le2$ ($|p|\le C/t$, $\hat C_t\le25t^2/4$, (A1)). With these inputs the proofs of \cref{app:a04:prop:AVa,app:a04:prop:AVb} go through verbatim, except for the weight mismatch of the pair: $f^{(t)}$ is $5t$-periodic with $5t$-block average $1$, $f^{(t/5)}$ is $t$-periodic with $t$-cell average $1$ \cite[(3.1)]{K2}. \emph{Weight normalisation.} For an integrand $F'$ on the finer lattice and $F$ on the coarser one,
\begin{gather*}
\textstyle\sum_{x'}(t/5)^3f'(x')F'(x')=\sum_{x\in\Lambda_t}t^3F'(x)+\sum_xt^35^{-3}\sum_{x'\in c(x)}f'(x')[F'(x')-F'(x)],\\
\textstyle\sum_xt^3f(x)F(x)=\sum_xt^3F(x)+\sum_{\text{blocks }B}\sum_{x\in B}t^3(f(x)-1)[F(x)-F(x_B)].
\end{gather*}
Both remainders are $\le\norm f_\infty$ times the in-cell, resp.\ in-block, oscillation of the integrand, which is step (4) of the proof of \cref{app:a04:prop:AVb} over $5t$-blocks: $O(t\cdot\mathrm{polylog})$ (near-diagonal points $\le Ct^2$). The same reduction applies to the $(Q^fG)$ lines. Hence the increment bound holds for $f_D$ with the same rate $t^{1/2}$.
\end{proof}

\begin{proof}[Proof of \cref{app:a04:lem:B2}]
By \cref{app:a04:lem:bulkform}, $\beta_{(2,0)}=(\beta_{(2,0)}+\mathfrak a)-\mathfrak a$, the first bracket being the original-vertex part. For it, \cref{app:a04:thm:B2orig} gives the sup bound $Ce^2\eta^{-1}(1+j)^{14}$ and the increment bound $C_\vartheta e^25^{-\vartheta j}\eta^{-1}$ for every $\vartheta\in(0,1)$, for both weights, for $1\le j\le k$. For the averaging part, \cref{app:a04:prop:AVa,app:a04:prop:AVb,app:a04:prop:AVf} give $\norm{\mathfrak a}_\infty\le Ce^2\eta^{-1}(1+j)^q$ and $|\bar{\mathfrak a}^{(j)}_k-\bar{\mathfrak a}^{(j+1)}_{k+1}|\le Ce^2\eta^{-1}[5^{-j/2}(1+j)^q+5^{-2j}]\le C_\vartheta e^25^{-\vartheta j}\eta^{-1}$ for every $\vartheta\in(0,\frac12)$, for both weights, for $1\le j\le k-J_0$. Adding, with $q$ the larger of $14$ and the exponent of the averaging part, gives both bounds of \cref{app:a04:lem:B2} with the constants depending on $\vartheta$, $\mathfrak a$, $e$, $m$, $\lmu$ and $m_0$ (and on $\norm{f_D}_\infty\le(5/3)^3$).
\end{proof}
